At this point the target binary has been modified with :
\begin{itemize}
    \item A DT\_NEEDED entry has been added to \textit{.dynamic} pointing to a shared object that contains 
        the forkserver declared as constructor, alongside a RUNPATH that points to that shared object.
    \item A symbol that represent the pointer to the \gls{SHM} has been inserted into \textit{.dynsym}
    \item A new section that will be used as personalized \gls{GOT} is added
    \item A relocation thank link the symbol to be resolved to the new section is added
\end{itemize}
Now the \gls{IP}s must be identified so that they can subsequently be instrumented.
First, a distinction must be made between two types of instrumentation : 
\begin{itemize}
    \item An initial instrumentation is needed in order to initialize some variables needed for the \gls{SHM} update rule.
    \item All the other instrumentations performs the actual \gls{SHM} update.
\end{itemize}

In the first case the \gls{IP} must be as close as possible to the \texttt{main} function of the binary.
Find the main address is easy both for the stripped and non stripped binaries, so, identified the main address, 
the tool tries to find the first address after the function preamble.
This goal is not achieved everytime because there is not a unified grammar for the function prologue. 
It depends on how many registers the callee wants to save on the stack, how many variables need the specific function \dots.
In general the first \gls{IP} lay in the first \gls{BB} of the main function.\\
All the other \gls{IP} are found with a with a \gls{DFS} starting from the initial main \gls{BB} and search between the 
successors \gls{BB} of the current one.
Specifically, once a \gls{BB} is found and is not been analyzed or discovered before, is added to a queue.
\begin{algorithm}[ht]
    \caption{Instrumentation point search with DFS algorithm}
    \label{alg:ip_search_dfs}
    \begin{algorithmic}[1]
        \STATE $found \gets[\,\,]$ \COMMENT {\small \textcolor{green!60!black}{found is a list of discovered basic blocks. It contains both the basic blocks to be instrumented and the basic block for which an instrumentation is not required}}
        \STATE $to\_instrument\gets[\,\,]$ \COMMENT{\small \textcolor{green!60!black}{this is the list of basic blocks to be instrumented. Is a subset of ``found'' list}}
        \STATE $current=\texttt{main\_bb}$
        \STATE $stack\gets[\,\,]$
        \STATE $successors\gets \textsc{ExtractSuccessors}(current,found)$
        \FOR{$s \in successors$}
            \STATE $stack.\textsc{push}((current,s))$
        \ENDFOR
        \WHILE{$stack \text{ is not empty}$}
            \STATE $current\_edge\gets stack.\textsc{pop}()$
            \STATE $prev \gets current\_edge[\,0\,] , curr \gets current\_edge[\,1\,]$
            \IF{$prev.last\_instruction.\_opcode \neq ``call''$}
                \STATE $to\_instrument.\textsc{append}(curr)$  \COMMENT{\small \textcolor{green!60!black}{If the previous basic block ends with a call instruction, the current basic block is not instrumented}}
            \ENDIF
            \STATE $found.\textsc{append}(curr)$
            \STATE $successors\gets \textsc{ExtractSuccessors}(current,found)$
            \FOR{$s \in successors$}
                \STATE $stack.\textsc{push}((current,s))$ 
            \ENDFOR
        \ENDWHILE
    \end{algorithmic}
\end{algorithm}
The \gls{IP}s discovery procedure could be seen in \cref{alg:ip_search_dfs}.\\
Two different lists must be used in order to separate between found basic blocks and basic blocks to be instrumented.
For example if a basic block ends with a call is ``useless'' to instrument the next one.
Infact a \texttt{call} is like a branch always taken.
If an edge that is always taken is instrumented, this would not add information to a fuzzer.
Infact if, during fuzzing, the first basic block is discovered, surely will be discovered also the second.
A difference could be with \texttt{jmp} instruction. Infact, although also in this case is a branch always taken, this could be 
useful for example in loops to disambiguate between the backward edge (return at the beginning of the loop) and the entry edge.
An example of this could be seen in \cref{fig:pow_example}.\\
The \textsc{ExtractSuccessors} function will be seen in the next chapter where the implementation details will be discussed 
in greater depth. 
For the moment, it can be said that this function filters the successors of the current basic block, 
discarding those that it would make no sense to consider.
For example, if in a certain point the binary performs a call to a library function (e.g. \texttt{printf}), 
without excluding this case, the program will instrument also the \textit{plt} stub of the library function.
\begin{figure}[t]
    \centering
    \begin{subfigure}[h]{0.8\textwidth}
        \centering
        \includegraphics[width=\textwidth]{pow_main_example.png}
        \caption{Main CFG}
        \label{fig:pow_main_example}
    \end{subfigure}
    \begin{subfigure}[h]{0.8\textwidth}
        \centering
        \includegraphics[width=\textwidth]{pow_function_example.png}
        \caption{Pow function CFG}
        \label{fig:pow_function_example}
    \end{subfigure}
    \caption{Example of the CFG of a program that compute the pow between two numbers}
    \label{fig:pow_example}
\end{figure}

In \cref{fig:pow_example} can be seen a disassembled \gls{CFG} obtained with Cutter \cite{cutter} that represent a program to compute the power operation.
With respect to \cref{fig:pow_main_example}, the \gls{IP}s found are: 
\begin{itemize}
    \item \textbf{\textit{0x11c2}} : initial main \gls{BB}. Here will be place the initial instrumentation
    \item \textbf{\textit{0x11f1}} (prev=\textit{0x11c2},next=\textit{0x11f1}) 
    \item \textbf{\textit{0x11db}} (prev=\textit{0x11c2},next=\textit{0x11db}) 
    \item \textbf{\textit{0x1253}} 
        \begin{itemize}
            \item (prev=\textit{0x11f1},next=\textit{0x1253})
            \item (prev=\textit{0x11db},next=\textit{0x1253}) 
        \end{itemize} 
    
\end{itemize}
Apart from the call to \texttt{my\_pow} function, all the other functions are exported from libc, and so does not 
need to be instrumented.

With respect to \cref{fig:pow_function_example}, the \gls{IP}s found are:
\begin{itemize}
    \item \textbf{\textit{0x11b5}} 
        \begin{itemize}
            \item (prev=\textit{0x1189},next=\textit{0x11b5}) 
            \item (prev=\textit{0x11a7},next=\textit{0x11b5}) 
        \end{itemize}
    \item \textbf{\textit{0x11a7}} (prev=\textit{0x11b5},next=\textit{0x11a7}) 
    \item \textbf{\textit{0x11bd}} (prev=\textit{0x11b5},next=\textit{0x11bd}) 
\end{itemize}
The initial \gls{BB} of the function is not instrumented.